% Header: General study / work / project template
%%%%%%%%%%%%%%%%%%%%%%%%%%%%%%%%%%%%%%%%%%%%%%%%%%%%%%%%%%%%%%%%%%%

\documentclass[11pt,a4paper]{scrartcl}

% Layout
\usepackage[margin=2.5cm]{geometry}
\usepackage[onehalfspacing]{setspace}

% General packages
\usepackage{graphicx}
\usepackage{enumitem}
\usepackage{tcolorbox}
\tcbuselibrary{breakable}
\usepackage[breaklinks=true,colorlinks=true,linkcolor=blue,urlcolor=blue,citecolor=blue]{hyperref}
\usepackage{amsmath,amsthm,amssymb,mathtools}
\usepackage{braket}
\usepackage{icomma}
\usepackage[T1]{fontenc}
\usepackage[utf8]{inputenc}

% Language: choose one
% \usepackage[ngerman]{babel}
\usepackage[english]{babel}

% Units / tables / floats / references
\usepackage[locale = UK,space-before-unit=true,per-mode = symbol]{siunitx}
\usepackage{booktabs,multirow}
\usepackage{placeins}
\usepackage[natbib,abbreviate=true,doi=false,style=numeric-comp,giveninits=true,sorting=none]{biblatex}
\usepackage{csquotes}

% Optional bibliography
% \addbibresource{references.bib}

% Graphics paths
\graphicspath{{Bilder/} {images/} {figures/} {chapters/}}

% Optional math shortcuts
\newcommand{\R}{\mathbb{R}}
\newcommand{\N}{\mathbb{N}}
\newcommand{\Q}{\mathbb{Q}}
%\newcommand{\set}[1]{\left\{ #1 \right\}}
\newcommand{\abs}[1]{\left| #1 \right|}
\newcommand{\norm}[1]{\left\lVert #1 \right\rVert}
\newcommand{\ol}[1]{\overline{#1}}

% Optional units
\DeclareSIUnit{\dBm}{dBm}

% Box styles
\newtcolorbox{calculationbox}[1][]{
    title=Calculation,
    breakable,
    #1
}

\newtcolorbox{solutionbox}[1][]{
    title=Solution,
    breakable,
    #1
}

\newtcolorbox{answerbox}[1][]{
    title=Answer,
    breakable,
    #1
}

\newtcolorbox{notebox}[1][]{
    title=Notes,
    breakable,
    #1
}

\newtcolorbox{sidecalcbox}[1][]{
    title=Side Calculation,
    breakable,
    #1
}

%%%%%%%%%%%%%%%%%%%%%%%%%%%%%%%%%%%%%%%%%%%%%%%%%%%%%%%%%%%%%%%%%%%
\begin{document}

\title{Theoretical Physics 2 - Blatt 04}
\author{Tobias Laser}
\date{\today}
\maketitle
\vfill
\thispagestyle{empty}

\pagenumbering{arabic}

% ================================================================
\paragraph{10. Measurement Probability}
Given the three-dimensional wavefunction
\begin{align*}
    \psi(x, y, z) = N \exp \left(-\left(\frac{\abs{x}}{2a} + \frac{\abs{y}}{2b} + \frac{z}{2c}\right)\right)
\end{align*}
where $a, b, c$ are positive real numbers.

\begin{enumerate}[label=(\alph*)]
    \item Compute the normalization constant $N$.
    \begin{calculationbox}
        $\abs{N}^2 \int dx dy dz \exp\left(-\left(\frac{\abs{x}}{2a} + \frac{\abs{y}}{2b} + \frac{z}{2c}\right)\right) = 1$
        \begin{calculationbox}
            $\int\limits_0^{+\infty} dx \exp\left(-\frac{x}{a}\right)$
        \end{calculationbox}
        $\abs{N}^2 2^3 a b c = 1 \Leftrightarrow N = \frac{1}{\sqrt{2^3 a b c}}$
    \end{calculationbox}
    \item What is the probability that a measurement of $X$ yields a result between $0$ and $a$?
    \begin{calculationbox}
        $\int\limits_0^a dx \frac{1}{2^3 a b c} \exp\left(\frac{-x}{a} \right) \int\limits dy dx \exp\left(-\left(\frac{\abs{y}}{b} + \frac{\abs{z}}{c} \right) \right) = \int\limits_0^a dx \frac{\exp\left(-\frac{x}{a} \right)}{2a} = -\frac{a \exp\left(-\frac{x}{a} \right)}{2 a}\left. \right|_0^a = \frac{1}{2}(1 - \frac{1}{e})$
    \end{calculationbox}
    \item What is the probability that a simultaneous measurement of $Y$ and $Z$ yields values between $-b$ and $b$, and between $-c$ and $c$, respectively?
    \begin{calculationbox}
        $\int\limits_{-\infty}^{+\infty} dx \int\limits_{-\infty}^{+\infty} dy \int\limits_{-\infty}^{+\infty} dz \frac{1}{2^3 a b c} \exp\left(-\left(\frac{\abs{x}}{a} + \frac{\abs{y}}{b} + \frac{\abs{z}}{c}\right) \right) = \frac{4}{4 b c} \int\limits_0^b dy \exp\left(-\frac{y}{b} \right) \int\limits_0^c dz \exp\left(-\frac{z}{c} \right) = \frac{1}{b c} \left(-b \exp\left(-\frac{y}{b} \right) \right)\left. \right|_0^b \left(-c \exp\left(-\frac{z}{c} \right) \right)\left. \right|_0^c = (1-\frac{1}{e})^2$
    \end{calculationbox}
    \item Compute the probability that a measurement of the momentum yields a result within the volume element $dp_x dp_y dp_z$, centered at the point $p_x = p_y = 0, p_z = \frac{\hbar}{c}$.
    \begin{calculationbox}
        Probability amplitude in momentum space:
        \begin{align*}
            \phi(\mathbf{p}) 
            &= \frac{1}{(2\pi\hbar)^{3/2}} \int dx\,dy\,dz \; \psi(x,y,z)\, e^{-i \frac{p_x x + p_y y + p_z z}{\hbar}}
        \end{align*}
        
        For $p_x = p_y = 0$, $p_z = \frac{\hbar}{c}$:
        \begin{align*}
            \phi(0,0,\hbar/c)
            &= \frac{1}{(2\pi\hbar)^{3/2}} \int dx\,dy\,dz \; \psi(x,y,z)\, e^{-i z / c}
        \end{align*}

        Insert $\psi(x,y,z) = N e^{-\left(\frac{|x|}{2a} + \frac{|y|}{2b} + \frac{|z|}{2c}\right)}$:
        \begin{align*}
            \phi 
            &= \frac{N}{(2\pi\hbar)^{3/2}} 
            \left( \int dx\, e^{-\frac{|x|}{2a}} \right)
            \left( \int dy\, e^{-\frac{|y|}{2b}} \right)
            \left( \int dz\, e^{-\frac{|z|}{2c}} e^{-i z / c} \right)
        \end{align*}

        Using
        \begin{align*}
            \int_{-\infty}^{\infty} e^{-\frac{|x|}{2a}} dx = 4a, 
            \quad 
            \int_{-\infty}^{\infty} e^{-\frac{|y|}{2b}} dy = 4b
        \end{align*}

        and splitting the $z$-integral:
        \begin{align*}
            \int_{-\infty}^{\infty} dz\, e^{-\frac{|z|}{2c}} e^{-i z / c}
            &= \int_{-\infty}^{0} dz\, e^{\frac{z}{2c}} e^{-i z / c}
            + \int_{0}^{\infty} dz\, e^{-\frac{z}{2c}} e^{-i z / c}
        \end{align*}

        \begin{align*}
            &= \int_{-\infty}^{0} dz\, e^{\left(\frac{1}{2c} - \frac{i}{c}\right) z}
            + \int_{0}^{\infty} dz\, e^{-\left(\frac{1}{2c} + \frac{i}{c}\right) z}
        \end{align*}

        Evaluating:
        \begin{align*}
            &= \frac{2c}{1 - 2i} + \frac{2c}{1 + 2i}
            = \frac{4c}{1 + 4}
            = \frac{4c}{5}
        \end{align*}

        Therefore:
        \begin{align*}
            \phi(0,0,\hbar/c)
            &= \frac{N}{(2\pi\hbar)^{3/2}} \cdot (4a)(4b)\cdot \frac{4c}{5}
            = \frac{64 abc\, N}{5 (2\pi\hbar)^{3/2}}
        \end{align*}

        Using $N = \frac{1}{\sqrt{8abc}}$:
        \begin{align*}
            \phi = \frac{64 abc}{5 (2\pi\hbar)^{3/2}} \cdot \frac{1}{\sqrt{8abc}}
            = \frac{64}{5 (2\pi\hbar)^{3/2}} \sqrt{\frac{abc}{8}}
        \end{align*}

        Probability:
        \begin{align*}
            P = |\phi|^2\, dp_x dp_y dp_z
            = \frac{64}{25} \frac{abc}{(\pi \hbar)^3} \, dp_x dp_y dp_z
        \end{align*}
    \end{calculationbox}
\end{enumerate}

% ================================================================
\paragraph{11. Operator Algebra and Measurements}
Consider a three-dimensional Hilbert space with an orthonormal basis $\{\phi_1,\phi_2,\phi_3\}$. The representations of two operators in this basis are
\begin{align*}
    H &= \hbar \omega_0
    \begin{pmatrix}
        1 & 0 & 0 \\
        0 & -1 & 0 \\
        0 & 0 & -1
    \end{pmatrix},
    &
    B &= b
    \begin{pmatrix}
        1 & 0 & 0 \\
        0 & 0 & 1 \\
        0 & 1 & 0
    \end{pmatrix}.
\end{align*}

\begin{enumerate}[label=(\alph*)]
    \item Compute the commutator of $H$ and $B$.
    \begin{calculationbox}
        We compute
        \begin{align*}
            [H,B] = HB - BH.
        \end{align*}
        Since
        \begin{align*}
            H &= \hbar\omega_0
            \begin{pmatrix}
                1 & 0 & 0 \\
                0 & -1 & 0 \\
                0 & 0 & -1
            \end{pmatrix},
            &
            B &= b
            \begin{pmatrix}
                1 & 0 & 0 \\
                0 & 0 & 1 \\
                0 & 1 & 0
            \end{pmatrix},
        \end{align*}
        we obtain
        \begin{align*}
            HB
            &= \hbar\omega_0 b
            \begin{pmatrix}
                1 & 0 & 0 \\
                0 & 0 & -1 \\
                0 & -1 & 0
            \end{pmatrix},
            \\
            BH
            &= \hbar\omega_0 b
            \begin{pmatrix}
                1 & 0 & 0 \\
                0 & 0 & -1 \\
                0 & -1 & 0
            \end{pmatrix}.
        \end{align*}
        Hence
        \begin{align*}
            [H,B] = 0.
        \end{align*}
    \end{calculationbox}

    \item Give the common eigenbasis of $H$ and $B$.
    \begin{calculationbox}
        The operator $H$ is already diagonal in the basis $\{\phi_1,\phi_2,\phi_3\}$:
        \begin{align*}
            H\phi_1 = \hbar\omega_0 \phi_1,
            \qquad
            H\phi_2 = -\hbar\omega_0 \phi_2,
            \qquad
            H\phi_3 = -\hbar\omega_0 \phi_3.
        \end{align*}

        For $B$, we have
        \begin{align*}
            B\phi_1 = b\phi_1,
        \end{align*}
        and on the subspace spanned by $\phi_2,\phi_3$,
        \begin{align*}
            B_0 =
            \begin{pmatrix}
                0 & 1 \\
                1 & 0
            \end{pmatrix}.
        \end{align*}
        Its eigenvalues are $\lambda = \pm 1$, with normalized eigenvectors
        \begin{align*}
            v_2 = \frac{1}{\sqrt{2}}(\phi_2 + \phi_3),
            \qquad
            v_3 = \frac{1}{\sqrt{2}}(\phi_2 - \phi_3).
        \end{align*}
        Also
        \begin{align*}
            v_1 = \phi_1.
        \end{align*}

        Therefore a common eigenbasis of $H$ and $B$ is
        \begin{align*}
            v_1 &= \phi_1,
            \\
            v_2 &= \frac{1}{\sqrt{2}}(\phi_2 + \phi_3),
            \\
            v_3 &= \frac{1}{\sqrt{2}}(\phi_2 - \phi_3).
        \end{align*}
    \end{calculationbox}

    \item Do the pairs $\{H,B\}$ and $\{H^2,B\}$ form a complete set of commuting observables?
    \begin{calculationbox}
        Since $[H,B]=0$ and both operators are Hermitian, they are compatible observables.

        For the pair $\{H,B\}$:
        \begin{align*}
            H v_1 &= \hbar\omega_0 v_1, & B v_1 &= b v_1,
            \\
            H v_2 &= -\hbar\omega_0 v_2, & B v_2 &= b v_2,
            \\
            H v_3 &= -\hbar\omega_0 v_3, & B v_3 &= -b v_3.
        \end{align*}
        The simultaneous eigenvalue pairs are therefore
        \begin{align*}
            (\hbar\omega_0,b), \qquad (-\hbar\omega_0,b), \qquad (-\hbar\omega_0,-b),
        \end{align*}
        which distinguish all basis vectors. Hence $\{H,B\}$ is a complete set of commuting observables.

        For the pair $\{H^2,B\}$:
        \begin{align*}
            H^2 = \hbar^2 \omega_0^2
            \begin{pmatrix}
                1 & 0 & 0 \\
                0 & 1 & 0 \\
                0 & 0 & 1
            \end{pmatrix}
            = \hbar^2 \omega_0^2 I.
        \end{align*}
        Thus $H^2$ has the same eigenvalue for all states and gives no additional information. Therefore $\{H^2,B\}$ is not complete.
    \end{calculationbox}

    \item The system is in the state
    \begin{align*}
        \psi = \left\{ \frac{1}{\sqrt{2}}, \frac{1}{2}, \frac{1}{2} \right\}.
    \end{align*}
    Compute the probabilities that:
    \begin{enumerate}[label=\roman*.]
        \item a measurement of $H$ gives the result $\hbar\omega_0$,
        \item a measurement of $H$ gives the result $-\hbar\omega_0$,
        \item a measurement of $B$ gives the result $b$,
        \item a measurement of $B$ gives the result $-b$,
        \item a measurement of $B$ gives the result $3b$.
    \end{enumerate}
    \begin{calculationbox}
        In the basis $\{\phi_1,\phi_2,\phi_3\}$,
        \begin{align*}
            \psi = \frac{1}{\sqrt{2}}\phi_1 + \frac{1}{2}\phi_2 + \frac{1}{2}\phi_3.
        \end{align*}

        For $H$:
        \begin{align*}
            P(H=\hbar\omega_0) &= \left|\frac{1}{\sqrt{2}}\right|^2 = \frac{1}{2},
            \\
            P(H=-\hbar\omega_0) &= \left|\frac{1}{2}\right|^2 + \left|\frac{1}{2}\right|^2 = \frac{1}{2}.
        \end{align*}

        For $B$, rewrite $\psi$ in the common eigenbasis:
        \begin{align*}
            \phi_2 &= \frac{1}{\sqrt{2}}(v_2 + v_3),
            \\
            \phi_3 &= \frac{1}{\sqrt{2}}(v_2 - v_3).
        \end{align*}
        Therefore
        \begin{align*}
            \psi
            &= \frac{1}{\sqrt{2}}v_1 + \frac{1}{2}\cdot\frac{1}{\sqrt{2}}(v_2+v_3)
            + \frac{1}{2}\cdot\frac{1}{\sqrt{2}}(v_2-v_3)
            \\
            &= \frac{1}{\sqrt{2}}v_1 + \frac{1}{\sqrt{2}}v_2.
        \end{align*}
        So
        \begin{align*}
            P(B=b) &= \left|\frac{1}{\sqrt{2}}\right|^2 + \left|\frac{1}{\sqrt{2}}\right|^2 = 1,
            \\
            P(B=-b) &= 0,
            \\
            P(B=3b) &= 0,
        \end{align*}
        since $3b$ is not an eigenvalue of $B$.
    \end{calculationbox}

    \item The system from (d) is measured twice.
    \begin{enumerate}[label=\roman*.]
        \item In the first measurement of $H$, one obtains $-\hbar\omega_0$. Compute the probability that a second measurement of $B$ gives $+b$.
        \item In the first measurement of $H$, one obtains $-\hbar\omega_0$. Compute the probability that a second measurement of $B$ gives $-b$.
        \item In the first measurement of $H$, one obtains $\hbar\omega_0$. What is obtained in the second measurement of $B$?
    \end{enumerate}
    \begin{calculationbox}
        If the first measurement of $H$ yields $-\hbar\omega_0$, the state collapses onto the normalized projection of $\psi$ onto the eigenspace spanned by $\phi_2,\phi_3$:
        \begin{align*}
            \psi' = \frac{\frac{1}{2}\phi_2 + \frac{1}{2}\phi_3}{\sqrt{\frac{1}{4}+\frac{1}{4}}}
            = \frac{1}{\sqrt{2}}(\phi_2+\phi_3)
            = v_2.
        \end{align*}
        Since $v_2$ is an eigenstate of $B$ with eigenvalue $+b$:
        \begin{align*}
            P(B=+b) = 1,
            \qquad
            P(B=-b) = 0.
        \end{align*}

        If the first measurement of $H$ yields $\hbar\omega_0$, the state collapses to
        \begin{align*}
            \psi'' = \phi_1 = v_1.
        \end{align*}
        Since $v_1$ is an eigenstate of $B$ with eigenvalue $+b$, the second measurement yields
        \begin{align*}
            B = +b
        \end{align*}
        with certainty.
    \end{calculationbox}

\end{enumerate}

% ================================================================
\paragraph{12. States of Minimal Position and Momentum Uncertainty}
The Heisenberg uncertainty relation for position and momentum operators can be proven as follows. For a given state $\ket{\psi}$ define
\begin{align*}
    \hat{X}' &= \hat{X} - \langle \hat{X} \rangle,
    &
    \hat{P}' &= \hat{P} - \langle \hat{P} \rangle,
\end{align*}
such that
\begin{align*}
    \Delta X^2 = \langle \hat{X}'^2 \rangle,
    \qquad
    \Delta P^2 = \langle \hat{P}'^2 \rangle.
\end{align*}
Consider the state
\begin{align*}
    \ket{\phi} = (\hat{X}' + i \lambda \hat{P}') \ket{\psi}
\end{align*}
with arbitrary real parameter $\lambda$.

\begin{enumerate}[label=(\alph*)]
    \item Show that from
    \begin{align*}
        \forall \lambda \in \mathbb{R}: \langle \phi | \phi \rangle \geq 0
    \end{align*}
    the Heisenberg uncertainty relation $\Delta X \Delta P \geq \hbar/2$ follows.
    \begin{calculationbox}
        We have
        \begin{align*}
            \langle \phi | \phi \rangle
            &= \bra{\psi}(\hat{X}' - i \lambda \hat{P}')(\hat{X}' + i \lambda \hat{P}')\ket{\psi}
            \geq 0.
        \end{align*}
        Expanding:
        \begin{align*}
            \langle \phi | \phi \rangle
            &= \bra{\psi}\hat{X}'^2\ket{\psi}
            + i\lambda \bra{\psi}\hat{X}'\hat{P}'\ket{\psi}
            - i\lambda \bra{\psi}\hat{P}'\hat{X}'\ket{\psi}
            + \lambda^2 \bra{\psi}\hat{P}'^2\ket{\psi}.
        \end{align*}
        Hence
        \begin{align*}
            \langle \phi | \phi \rangle
            &= \Delta X^2 + \lambda^2 \Delta P^2 + i\lambda \langle [\hat{X}',\hat{P}'] \rangle.
        \end{align*}
        Since constants commute,
        \begin{align*}
            [\hat{X}',\hat{P}'] = [\hat{X},\hat{P}] = i\hbar,
        \end{align*}
        therefore
        \begin{align*}
            \langle \phi | \phi \rangle
            &= \Delta X^2 + \lambda^2 \Delta P^2 + i\lambda (i\hbar)
            \\
            &= \Delta X^2 + \lambda^2 \Delta P^2 - \lambda \hbar.
        \end{align*}
        Thus for all $\lambda \in \mathbb{R}$,
        \begin{align*}
            \Delta X^2 + \lambda^2 \Delta P^2 - \lambda \hbar \geq 0.
        \end{align*}
        This quadratic expression in $\lambda$ must be non-negative for all real $\lambda$, so its discriminant must satisfy
        \begin{align*}
            \hbar^2 - 4 \Delta X^2 \Delta P^2 \leq 0.
        \end{align*}
        Hence
        \begin{align*}
            4 \Delta X^2 \Delta P^2 \geq \hbar^2,
        \end{align*}
        and therefore
        \begin{align*}
            \Delta X \Delta P \geq \frac{\hbar}{2}.
        \end{align*}
    \end{calculationbox}

    \item Show that equality holds for a state $\ket{\psi}$ if and only if
    \begin{align*}
        (\hat{X}' + i\lambda_0 \hat{P}')\ket{\psi} = 0
    \end{align*}
    with
    \begin{align*}
        \lambda_0 = \frac{\hbar}{2\Delta P^2} = \frac{2\Delta X^2}{\hbar}.
    \end{align*}
    \begin{calculationbox}
        Equality in the uncertainty relation holds if and only if the quadratic expression
        \begin{align*}
            f(\lambda) = \Delta X^2 + \lambda^2 \Delta P^2 - \lambda \hbar
        \end{align*}
        has a double root, i.e. its minimum is zero.

        The minimum is attained at
        \begin{align*}
            \frac{d f}{d\lambda} = 2\lambda \Delta P^2 - \hbar = 0
            \qquad \Longrightarrow \qquad
            \lambda_0 = \frac{\hbar}{2\Delta P^2}.
        \end{align*}
        At equality we also have
        \begin{align*}
            \Delta X \Delta P = \frac{\hbar}{2},
        \end{align*}
        so that
        \begin{align*}
            \Delta X^2 = \frac{\hbar^2}{4\Delta P^2}.
        \end{align*}
        Hence
        \begin{align*}
            \lambda_0 = \frac{\hbar}{2\Delta P^2} = \frac{2\Delta X^2}{\hbar}.
        \end{align*}

        Now, if
        \begin{align*}
            (\hat{X}' + i\lambda_0 \hat{P}')\ket{\psi} = 0,
        \end{align*}
        then $\ket{\phi}=0$ and therefore
        \begin{align*}
            \langle \phi | \phi \rangle = 0,
        \end{align*}
        so the lower bound is saturated.

        Conversely, if equality holds, then the minimum value of $\langle \phi|\phi\rangle$ is zero at $\lambda=\lambda_0$, which implies
        \begin{align*}
            \|(\hat{X}' + i\lambda_0 \hat{P}')\ket{\psi}\|^2 = 0.
        \end{align*}
        Therefore
        \begin{align*}
            (\hat{X}' + i\lambda_0 \hat{P}')\ket{\psi} = 0.
        \end{align*}
    \end{calculationbox}

    \item Write the equation
    \begin{align*}
        (\hat{X}' + i\lambda_0 \hat{P}')\ket{\psi} = 0
    \end{align*}
    in the position representation and solve the resulting differential equation for $\psi(x)=\langle x|\psi\rangle$. Consider the ansatz
    \begin{align*}
        \psi(x)=\exp\!\left(i\frac{\langle \hat{P}\rangle x}{\hbar}\right) f(x-\langle \hat{X}\rangle).
    \end{align*}
    Which states are obtained?
    \begin{calculationbox}
        In the position representation,
        \begin{align*}
            \hat{X}\psi(x) &= x\psi(x),
            &
            \hat{P}\psi(x) &= -i\hbar \frac{d}{dx}\psi(x).
        \end{align*}
        Therefore
        \begin{align*}
            \hat{X}' &= x - \langle \hat{X} \rangle,
            &
            \hat{P}' &= -i\hbar \frac{d}{dx} - \langle \hat{P} \rangle.
        \end{align*}
        The equation becomes
        \begin{align*}
            \left[
                x - \langle \hat{X} \rangle
                + i\lambda_0\left(-i\hbar\frac{d}{dx} - \langle \hat{P} \rangle\right)
            \right]\psi(x)=0.
        \end{align*}
        Hence
        \begin{align*}
            \left[
                x - \langle \hat{X} \rangle
                + \lambda_0 \hbar \frac{d}{dx}
                - i\lambda_0 \langle \hat{P} \rangle
            \right]\psi(x)=0.
        \end{align*}

        Use the ansatz
        \begin{align*}
            \psi(x)=\exp\!\left(i\frac{\langle \hat{P}\rangle x}{\hbar}\right) f(x-\langle \hat{X}\rangle).
        \end{align*}
        Let
        \begin{align*}
            \xi = x-\langle \hat{X}\rangle.
        \end{align*}
        Then
        \begin{align*}
            \frac{d\psi}{dx}
            = \exp\!\left(i\frac{\langle \hat{P}\rangle x}{\hbar}\right)
            \left(
                i\frac{\langle \hat{P}\rangle}{\hbar} f(\xi) + f'(\xi)
            \right).
        \end{align*}
        Substituting into the differential equation, the terms containing $\langle \hat{P}\rangle$ cancel, and we obtain
        \begin{align*}
            \xi f(\xi) + \lambda_0 \hbar f'(\xi) = 0.
        \end{align*}
        Thus
        \begin{align*}
            \frac{f'(\xi)}{f(\xi)} = - \frac{\xi}{\lambda_0 \hbar}.
        \end{align*}
        Integrating:
        \begin{align*}
            \ln f(\xi) = -\frac{\xi^2}{2\lambda_0 \hbar} + \text{const}.
        \end{align*}
        Therefore
        \begin{align*}
            f(\xi) = C \exp\!\left(-\frac{\xi^2}{2\lambda_0 \hbar}\right).
        \end{align*}
        Since
        \begin{align*}
            \lambda_0 = \frac{2\Delta X^2}{\hbar},
        \end{align*}
        we get
        \begin{align*}
            f(\xi) = C \exp\!\left(-\frac{\xi^2}{4\Delta X^2}\right).
        \end{align*}
        Hence
        \begin{align*}
            \psi(x)
            = C \exp\!\left(
                -\frac{(x-\langle \hat{X}\rangle)^2}{4\Delta X^2}
                + i\frac{\langle \hat{P}\rangle x}{\hbar}
            \right).
        \end{align*}

        To normalize:
        \begin{align*}
            1
            &= |C|^2 \int_{-\infty}^{\infty}
            \exp\!\left(-\frac{(x-\langle \hat{X}\rangle)^2}{2\Delta X^2}\right)\,dx
            \\
            &= |C|^2 \sqrt{2\pi}\,\Delta X.
        \end{align*}
        Thus
        \begin{align*}
            C = \frac{1}{(2\pi \Delta X^2)^{1/4}}
        \end{align*}
        up to an irrelevant overall phase.

        Therefore the states are
        \begin{align*}
            \psi(x)
            = \frac{1}{(2\pi \Delta X^2)^{1/4}}
            \exp\!\left(
                -\frac{(x-\langle \hat{X}\rangle)^2}{4\Delta X^2}
                + i\frac{\langle \hat{P}\rangle x}{\hbar}
            \right),
        \end{align*}
        i.e. Gaussian wave packets, also called minimum-uncertainty states.
    \end{calculationbox}
\end{enumerate}
\end{document}